\documentclass[../../../include/open-logic-section]{subfiles}

\begin{document}

\olfileid[pt]{his}{set}{pathology}
\olsection{Patologias}
% Interlinhado compacto: \PatologiaCompactSpread em open-logic-config.sty.
\begin{patologiascompactas}

Contudo, a definição de um \emph{limite} acabou por permitir construções bastante ``patológicas''.
\vspace{-1pt}%

Por volta da década de 1830, Bolzano descobriu uma função que era
\emph{contínua em todo o lado}, mas \emph{não diferenciável em parte nenhuma}.
(Infelizmente, Bolzano nunca publicou isto; a ideia foi pela primeira vez
encontrada pelos matemáticos em 1872, graças à descoberta independente
de Weierstrass da mesma ideia.)\footnote{A história está
documentada em notas de rodapé extremamente minuciosas do artigo da Wikipédia sobre a
\href{http://en.wikipedia.org/wiki/Weierstrass_function}{função
de Weierstrass}.} Isto era, no mínimo, bastante surpreendente.
É fácil encontrar funções, como $|x|$, que são contínuas
em todo o lado mas não diferenciáveis num dado ponto. Mas uma
função contínua em todo o lado e \emph{não diferenciável em parte nenhuma} é uma criatura muito diferente. Considere, por um momento, como
tentaria desenhar tal função. Para garantir que é contínua, deve
consegui-la desenhar sem nunca levantar a caneta da página;
mas para garantir que não é diferenciável em nenhum ponto, teria de mudar
abruptamente a direção da caneta, constantemente.
\vspace{-1pt}%

Seguiram-se mais ``patologias''. A 5 de janeiro de 1874, Cantor escreveu uma
carta a Dedekind, colocando o problema:
\vspace{-2pt}%
\begin{quote}
Poderá uma superfície (digamos um quadrado incluindo o seu contorno) estar correlacionada biunivocamente com uma linha (digamos uma reta incluindo as extremidades) de modo que a cada ponto da superfície corresponda um ponto da
linha e, reciprocamente, a cada ponto da linha corresponda um
ponto da superfície?

Ainda me parece, de momento, que a resposta a esta questão é
muito difícil---embora aqui também se esteja tão impelido a dizer \emph{não}
que se queira considerar quase supérflua a prova.
[Citado em \citealt{Gouvea2011}]
\end{quote}
% Menos compressão que -4pt: alinha o salto quote→corpo ao ar dos outros parágrafos.
\vspace{-2pt}%
Mas, em 1877, Cantor provou que estivera errado. De fato, uma linha
e um quadrado têm exatamente o mesmo número de pontos. Escreveu a 29
de junho de 1877 a Dedekind ``\emph{je le vois, mais je ne le crois pas}'';
isto é, ``eu vejo, mas não acredito''. Na ``história recebida''
da matemática, isto é frequentemente tomado como indicador de quão
\emph{literalmente incríveis} eram estes novos resultados para os
matemáticos da época. (A correspondência é apresentada em
\citet{Gouvea2011}, e voltaremos a ela em
\olref[his][set][mythology]{sec}. A prova de Cantor está esboçada em
\olref[his][set][cantorplane]{sec}.)
% Espaço local refs (Gouvêa / secções)→«Inspirado…»; ~parskip do livro; compensado na grelha.
\vspace{0.22\baselineskip}%

Inspirado pelo resultado de Cantor, Peano começou a considerar se
seria possível mapear suavemente uma linha sobre um plano. Isto
seria uma \emph{curva que preenche o espaço}. Em \citeyear{Peano1890},
Peano construiu precisamente tal curva. Isto é verdadeiramente contraintuitivo:
Euclides tinha definido uma linha como ``comprimento sem largura'' (Livro I,
Definição 2), mas Peano mostrara que, enrolando adequadamente uma linha,
o seu comprimento pode transformar-se em largura. Em
\citeyear{Hilbert1891}, Hilbert descreveu uma curva preenchedora de espaço ligeiramente mais intuitiva,
juntamente com algumas figuras que a ilustram. A
curva é construída em sequência, e eis as primeiras seis fases da
construção:

\end{patologiascompactas}
% Folga na grelha: compensa quote→corpo (-2pt em vez de -4pt) e 0,22\baselineskip antes de «Inspirado».
\enlargethispage{\dimexpr 3\baselineskip+7pt\relax}%
\begin{center}
\begin{tikzpicture}[scale=0.80]
\begin{scope}[xshift=20pt, yshift=20pt]
	\draw [oldiagcolorC, l-system={Hilbert curve, step=40pt, angle=90, axiom=L, order=1}] lindenmayer system;
\end{scope}
\begin{scope}[xshift=110pt, yshift=10pt]
	\draw [oldiagcolorC, l-system={Hilbert curve, step=20pt, angle=90, axiom=L, order=2}] lindenmayer system;
\end{scope}
\begin{scope}[xshift=205pt, yshift=5pt]
	\draw [oldiagcolorC, l-system={Hilbert curve, step=10pt, angle=90, axiom=L, order=3}] lindenmayer system;
\end{scope}
\begin{scope}[xshift=2.5pt, yshift=-97.5pt]
	\draw [oldiagcolorC, l-system={Hilbert curve, step=5pt, angle=90, axiom=L, order=4}] lindenmayer system;
\end{scope}
\begin{scope}[xshift=101.25pt, yshift=-98.75pt]
	\draw [oldiagcolorC, l-system={Hilbert curve, step=2.5pt, angle=90, axiom=L, order=5}] lindenmayer system;
\end{scope}
\begin{scope}[xshift=200.625pt, yshift=-99.375pt]
	\draw [oldiagcolorC, l-system={Hilbert curve, step=1.25pt, angle=90, axiom=L, order=6}] lindenmayer system;
\end{scope}
\draw[gray] (0pt,0pt) rectangle (80pt, 80pt);
\draw[gray] (100pt,0pt) rectangle (180pt, 80pt);
\draw[gray] (200pt,0pt) rectangle (280pt, 80pt);
\draw[gray] (0pt,-100pt) rectangle (80pt, -20pt);
\draw[gray] (100pt,-100pt) rectangle (180pt, -20pt);
\draw[gray] (200pt,-100pt) rectangle (280pt, -20pt);
\end{tikzpicture}
\end{center}

\noindent
No limite---uma noção que, entretanto, recebera definição rigorosa---todo o quadrado fica inteiramente preenchido na cor !!{colorC}. E, de passagem, a curva de Hilbert é contínua em todo o lado mas não diferenciável em
parte nenhuma; intuitivamente porque, no limite infinito, a função
muda abruptamente de direção a cada instante. (Esboçaremos a construção de Hilbert
com mais pormenor na \olref[his][set][hilbertcurve]{sec}.)

\enlargethispage{2\baselineskip}%
% PT mais longo que EN: se o bloco final não couber, passa inteiro à página
% seguinte (evita viúva «propaganda» / linha órfã no topo da pág. seguinte).
\needspace{8\baselineskip}%
Para o bem ou para o mal, estas construções geométricas ``patológicas''
foram tratadas como motivo para duvidar de apelos à intuição geométrica. Tornaram-se
algo próximo de \emph{propaganda} por uma nova forma de fazer
matemática, que culminaria na teoria dos conjuntos. Na construção posterior do mito em torno do assunto, repetia-se, frequentemente, que estes
resultados eram ao mesmo tempo perfeitamente rigorosos e perfeitamente chocantes. Serviam
portanto um duplo propósito: como aviso contra confiar na
intuição geométrica, e como uma demonstração da fertilidade de novas
formas de pensar.

\end{document}
